\documentclass[12pt,a4paper]{article}
\usepackage[utf8]{inputenc}
\usepackage{geometry}
\usepackage{graphicx}
\usepackage{listings}
\usepackage{xcolor}
\usepackage{hyperref}
\usepackage{enumitem}
\usepackage{titlesec}
\usepackage{amsmath}
\usepackage{tcolorbox}

\geometry{margin=1in}

% Code listing style
\lstset{
    language=bash,
    basicstyle=\ttfamily\small,
    keywordstyle=\color{blue},
    commentstyle=\color{green!60!black},
    stringstyle=\color{red},
    numbers=left,
    numberstyle=\tiny\color{gray},
    stepnumber=1,
    numbersep=5pt,
    backgroundcolor=\color{gray!10},
    frame=single,
    breaklines=true,
    breakatwhitespace=true,
    tabsize=2,
    showstringspaces=false
}

% Python code style
\lstdefinestyle{python}{
    language=Python,
    basicstyle=\ttfamily\small,
    keywordstyle=\color{blue},
    commentstyle=\color{green!60!black},
    stringstyle=\color{red},
    numbers=left,
    numberstyle=\tiny\color{gray},
    backgroundcolor=\color{gray!10},
    frame=single,
    breaklines=true
}

% Custom boxes for tips and warnings
\newtcolorbox{tipbox}{
    colback=blue!5!white,
    colframe=blue!75!black,
    title=\textbf{💡 Tip},
    fonttitle=\bfseries
}

\newtcolorbox{warningbox}{
    colback=orange!5!white,
    colframe=orange!75!black,
    title=\textbf{⚠️ Warning},
    fonttitle=\bfseries
}

\newtcolorbox{infobox}{
    colback=green!5!white,
    colframe=green!75!black,
    title=\textbf{ℹ️ Information},
    fonttitle=\bfseries
}

\title{\Huge\textbf{Complete Guide: Running Hand Gesture Recognition\\on Raspberry Pi}}
\author{Step-by-Step Beginner's Guide}
\date{\today}

\begin{document}

\maketitle

\tableofcontents
\newpage

\section{Introduction}

\subsection{What is This Guide About?}

Imagine you have a magic camera that can understand hand gestures! This guide will help you set up a hand gesture recognition system on a Raspberry Pi (a small, affordable computer). Think of it like teaching your computer to recognize when you wave, point, or make different hand shapes.

\begin{infobox}
\textbf{What You'll Learn:}
\begin{itemize}
    \item How to prepare your Raspberry Pi
    \item How to install all the necessary software
    \item How to transfer your project files
    \item How to run the gesture recognition program
    \item How to fix common problems
\end{itemize}
\end{infobox}

\subsection{What Do You Need?}

Before we start, make sure you have:

\begin{itemize}
    \item \textbf{Raspberry Pi 4 or 5} (with power supply)
    \item \textbf{MicroSD card} (32GB or larger, with Raspberry Pi OS installed)
    \item \textbf{USB webcam} or Raspberry Pi Camera Module
    \item \textbf{Keyboard and mouse} (USB)
    \item \textbf{Monitor} (with HDMI cable)
    \item \textbf{Internet connection} (WiFi or Ethernet cable)
    \item \textbf{Your project files} (from your computer)
\end{itemize}

\begin{tipbox}
If you don't have a monitor, you can use SSH (remote connection) from your computer to control the Raspberry Pi. We'll explain this too!
\end{tipbox}

\section{Understanding the Project}

\subsection{What Does This Project Do?}

This project uses something called \textbf{MediaPipe} (a special tool made by Google) to:
\begin{enumerate}
    \item Look at your hand through a camera
    \item Find important points on your hand (like fingertips)
    \item Understand what gesture you're making
    \item Show the result on the screen
\end{enumerate}

Think of it like this: The camera is the "eyes," MediaPipe is the "brain" that understands hands, and the program shows you what it sees!

\subsection{What Files Do We Need?}

Your project has these important parts:

\begin{verbatim}
project_folder/
├── app.py                    ← The main program (this is what we'll run!)
├── model/                    ← The "brain" files
│   ├── keypoint_classifier/
│   │   ├── keypoint_classifier.tflite
│   │   └── keypoint_classifier_label.csv
│   └── point_history_classifier/
│       ├── point_history_classifier.tflite
│       └── point_history_classifier_label.csv
└── utils/                    ← Helper tools
    └── cvfpscalc.py
\end{verbatim}

\begin{warningbox}
\textbf{Important:} Make sure you have ALL these files! The program won't work without them.
\end{warningbox}

\section{Step 1: First Time Setup of Raspberry Pi}

\subsection{What is This Step?}

This is like turning on a new computer for the first time. We need to make sure everything is working!

\subsection{Step 1.0: Two Ways to Use Your Raspberry Pi}

Before we start, you should know there are two main ways to control your Raspberry Pi:

\subsubsection{Method A: Desktop Mode (With Monitor)}

If you have a monitor, keyboard, and mouse connected directly to your Pi:

\begin{enumerate}
    \item You'll see a desktop (like Windows or Mac)
    \item Click the Terminal icon (black box with \texttt{>\_}) in the top menu bar
    \item Or press \texttt{Ctrl + Alt + T} to open Terminal
    \item You'll see: \texttt{pi@raspberrypi:\~\$}
\end{enumerate}

\begin{infobox}
\textbf{Desktop Mode Benefits:}
\begin{itemize}
    \item See everything visually
    \item Easy to use if you're new to Linux
    \item Can see the camera window directly
\end{itemize}
\end{infobox}

\subsubsection{Method B: SSH (Remote from Your Computer)}

If you want to control your Pi from your Windows/Mac computer (no monitor needed):

\textbf{Step 1: Find Your Pi's IP Address}

On your Pi (if you have access), open Terminal and type:
\begin{lstlisting}
hostname -I
\end{lstlisting}

You'll see something like: \texttt{192.168.1.100}

\textbf{Step 2: Enable SSH on Pi}

If SSH is not enabled:
\begin{lstlisting}
sudo systemctl enable ssh
sudo systemctl start ssh
\end{lstlisting}

\textbf{Step 3: Connect from Your Computer}

\textbf{On Windows:}
\begin{enumerate}
    \item Open PowerShell or Command Prompt
    \item Type:
    \begin{lstlisting}
    ssh pi@192.168.1.100
    \end{lstlisting}
    (Replace with your Pi's actual IP address)
    \item Enter password when asked (default is \texttt{raspberry})
\end{enumerate}

\textbf{On Mac/Linux:}
\begin{lstlisting}
ssh pi@192.168.1.100
\end{lstlisting}

\begin{tipbox}
Once connected via SSH, you'll see the same terminal prompt and can run all the same commands!
\end{tipbox}

\subsection{Step 1.1: Turn On Your Raspberry Pi}

\begin{enumerate}
    \item Connect your Raspberry Pi to:
    \begin{itemize}
        \item Power supply (plug it in!)
        \item Monitor (using HDMI cable)
        \item Keyboard and mouse (USB ports)
        \item Internet (WiFi or Ethernet cable)
    \end{itemize}
    
    \item Turn it on by plugging in the power supply
    \item Wait for it to start (you'll see a desktop appear)
\end{enumerate}

\begin{tipbox}
If nothing appears on the screen, check:
\begin{itemize}
    \item Is the power LED on? (should be red)
    \item Is the HDMI cable connected properly?
    \item Is the monitor turned on?
\end{itemize}
\end{tipbox}

\subsection{Step 1.2: Open the Terminal}

The terminal is like a command center where we type instructions to the computer.

\begin{enumerate}
    \item Look for a black icon with a \texttt{>\_} symbol (usually in the top menu bar)
    \item Click on it to open the Terminal
    \item You should see something like: \texttt{pi@raspberrypi:\~\$}
\end{enumerate}

\begin{infobox}
\textbf{What you're seeing:}
\begin{itemize}
    \item \texttt{pi} = your username
    \item \texttt{raspberrypi} = your computer's name
    \item \texttt{\~} = you're in your home folder
    \item \texttt{\$} = ready for your command
\end{itemize}
\end{infobox}

\subsection{Step 1.2.1: Essential Linux Commands You'll Need}

Before we continue, here are the most important commands you'll use. Think of these as your toolkit:

\textbf{Navigation Commands:}
\begin{lstlisting}
pwd                    # Shows where you are (current folder)
ls                     # List files in current folder
ls -la                 # List all files (including hidden ones)
cd folder_name         # Go into a folder
cd ..                  # Go up one folder level
cd ~                   # Go to your home folder
cd /home/pi            # Go to a specific location
\end{lstlisting}

\textbf{File Operations:}
\begin{lstlisting}
mkdir folder_name      # Create a new folder
rm file_name           # Delete a file
rm -r folder_name      # Delete a folder (be careful!)
cp file1 file2         # Copy a file
mv file1 file2         # Move or rename a file
nano file.txt          # Open text editor to edit file
cat file.txt           # View file contents
\end{lstlisting}

\textbf{System Commands:}
\begin{lstlisting}
sudo command           # Run as administrator (needs password)
sudo apt update        # Update list of available packages
sudo apt upgrade       # Upgrade installed packages
sudo reboot            # Restart Raspberry Pi
sudo shutdown now      # Shut down Raspberry Pi
\end{lstlisting}

\textbf{Getting Help:}
\begin{lstlisting}
command --help         # Get help for any command
man command            # Read the manual page
\end{lstlisting}

\begin{tipbox}
\textbf{Remember:}
\begin{itemize}
    \item Press \texttt{Tab} to auto-complete file/folder names
    \item Press \texttt{Up Arrow} to see previous commands
    \item Press \texttt{Ctrl+C} to stop a running program
    \item Press \texttt{Ctrl+D} to exit terminal
\end{itemize}
\end{tipbox}

\subsection{Step 1.3: Check Your Python Version}

Python is the programming language we'll use. Let's check if it's installed:

\begin{lstlisting}
python3 --version
\end{lstlisting}

You should see something like: \texttt{Python 3.9.2} or \texttt{Python 3.11.5}

\begin{warningbox}
\textbf{Important:} MediaPipe works best with Python 3.9, 3.10, or 3.11. If you have Python 3.12 or higher, you might have problems. If you see a version that's too new, we'll fix it later!
\end{warningbox}

\section{Step 2: Update Your Raspberry Pi}

\subsection{What is This Step?}

Think of this like updating your phone's apps. We're making sure everything is the latest version so things work smoothly.

\subsection{Step 2.1: Update the System}

Type these commands one by one (press Enter after each):

\begin{lstlisting}
sudo apt update
\end{lstlisting}

This downloads a list of available updates. It might ask for your password (the default is usually \texttt{raspberry}).

\begin{lstlisting}
sudo apt upgrade -y
\end{lstlisting}

This actually installs the updates. The \texttt{-y} means "yes, do it automatically."

\begin{tipbox}
This might take 10-30 minutes. It's like downloading and installing many app updates at once. Be patient! You can see progress bars moving.
\end{tipbox}

\section{Step 3: Install System Dependencies}

\subsection{What is This Step?}

Dependencies are like tools and helpers that our program needs to work. It's like needing a hammer and nails to build something - you can't build without the right tools!

\subsection{Step 3.1: Install Basic Tools}

Type this command:

\begin{lstlisting}
sudo apt install -y build-essential cmake python3-pip python3-venv
\end{lstlisting}

\begin{infobox}
\textbf{What we're installing:}
\begin{itemize}
    \item \texttt{build-essential} = tools to build/compile software
    \item \texttt{cmake} = a tool that helps build software
    \item \texttt{python3-pip} = a tool to install Python packages
    \item \texttt{python3-venv} = tool to create isolated Python environments
\end{itemize}
\end{infobox}

\subsection{Step 3.2: Install OpenCV Dependencies}

OpenCV is like special glasses that help the computer "see" images from the camera.

\begin{lstlisting}
sudo apt install -y libopencv-dev python3-opencv libatlas-base-dev
\end{lstlisting}

\subsection{Step 3.3: Install Additional Libraries}

These are extra helpers for handling images and videos:

\begin{lstlisting}
sudo apt install -y libjpeg-dev libtiff5-dev libjasper-dev libpng12-dev
sudo apt install -y libavcodec-dev libavformat-dev libswscale-dev libv4l-dev
sudo apt install -y libxvidcore-dev libx264-dev
\end{lstlisting}

\begin{tipbox}
You can copy and paste these commands. Just make sure to press Enter after each one!
\end{tipbox}

\section{Step 4: Create a Virtual Environment}

\subsection{What is This Step?}

A virtual environment is like a separate room where we keep all our project's tools. This way, they don't mix with other projects' tools!

\subsection{Step 4.1: Create a Project Folder}

First, let's make a special folder for our project:

\begin{lstlisting}
mkdir ~/gesture_recognition
cd ~/gesture_recognition
\end{lstlisting}

\begin{infobox}
\textbf{What we did:}
\begin{itemize}
    \item \texttt{mkdir} = "make directory" (create a folder)
    \item \texttt{\~} = your home folder (shortcut)
    \item \texttt{cd} = "change directory" (go into the folder)
\end{itemize}
\end{infobox}

\subsection{Step 4.2: Create the Virtual Environment}

Now let's create our special "room":

\begin{lstlisting}
python3 -m venv venv_mediapipe_pi
\end{lstlisting}

This creates a folder called \texttt{venv\_mediapipe\_pi} with a clean Python environment inside.

\begin{tipbox}
This might take a minute. The computer is setting up a fresh Python installation just for this project!
\end{tipbox}

\subsection{Step 4.3: Activate the Virtual Environment}

Now let's "enter" this special room:

\begin{lstlisting}
source venv_mediapipe_pi/bin/activate
\end{lstlisting}

You should see \texttt{(venv\_mediapipe\_pi)} appear at the beginning of your command line. This means you're "inside" the virtual environment!

\begin{warningbox}
\textbf{Important:} Every time you open a new terminal, you need to run this command again to activate the virtual environment!
\end{warningbox}

\section{Step 5: Install Python Packages}

\subsection{What is This Step?}

Now we're installing the actual Python tools (called "packages") that our program needs. Think of them as apps for Python!

\subsection{Step 5.1: Upgrade pip}

First, let's make sure our package installer is up to date:

\begin{lstlisting}
pip install --upgrade pip setuptools wheel
\end{lstlisting}

\subsection{Step 5.2: Install NumPy}

NumPy helps with math and arrays (like lists of numbers):

\begin{lstlisting}
pip install numpy
\end{lstlisting}

\begin{tipbox}
This might take 2-5 minutes. NumPy is being compiled (built) for your Raspberry Pi!
\end{tipbox}

\subsection{Step 5.3: Install OpenCV}

OpenCV helps the computer see and process images:

\begin{lstlisting}
pip install opencv-python-headless
\end{lstlisting}

We use "headless" version because Raspberry Pi doesn't need the GUI parts (we're running it without a desktop display sometimes).

\subsection{Step 5.4: Install TensorFlow Lite}

TensorFlow Lite is a lightweight version of TensorFlow (a machine learning tool):

\begin{lstlisting}
pip install tensorflow-lite-runtime
\end{lstlisting}

\begin{infobox}
We use the "Lite" version because it's much smaller and faster on Raspberry Pi. The full TensorFlow would be too heavy!
\end{infobox}

\subsection{Step 5.5: Install MediaPipe}

This is the big one! MediaPipe is what recognizes your hand gestures:

\begin{lstlisting}
pip install mediapipe
\end{lstlisting}

\begin{warningbox}
\textbf{This will take 15-30 minutes!} MediaPipe is a large package and needs to be compiled. Be very patient. You'll see lots of text scrolling - that's normal!
\end{warningbox}

\begin{tipbox}
If this fails, try:
\begin{lstlisting}
pip install mediapipe==0.10.8
\end{lstlisting}
This installs a specific version that's known to work well.
\end{tipbox}

\subsection{Step 5.6: Verify Installation}

Let's check if everything installed correctly:

\begin{lstlisting}
python3 -c "import cv2; print('OpenCV:', cv2.__version__)"
python3 -c "import numpy; print('NumPy:', numpy.__version__)"
python3 -c "import mediapipe as mp; print('MediaPipe:', mp.__version__)"
\end{lstlisting}

If you see version numbers, everything is good! If you see errors, we'll fix them in the troubleshooting section.

\section{Step 6: Transfer Your Project Files}

\subsection{What is This Step?}

Now we need to copy your project files from your computer to the Raspberry Pi. It's like moving your homework from one computer to another!

\subsection{Step 6.1: Prepare Files on Your Computer}

On your Windows/Mac computer, make sure you have these folders ready:

\begin{verbatim}
Your Project/
├── app.py
├── model/
│   ├── __init__.py
│   ├── keypoint_classifier/
│   │   ├── keypoint_classifier.tflite
│   │   ├── keypoint_classifier_label.csv
│   │   └── keypoint_classifier.py
│   └── point_history_classifier/
│       ├── point_history_classifier.tflite
│       ├── point_history_classifier_label.csv
│       └── point_history_classifier.py
└── utils/
    ├── __init__.py
    └── cvfpscalc.py
\end{verbatim}

\begin{warningbox}
\textbf{Don't copy the venv folders!} They're too big and we'll create a new one on the Pi.
\end{warningbox}

\subsection{Step 6.2: Method 1 - Using USB Drive}

This is the easiest method:

\begin{enumerate}
    \item Copy your project folder to a USB drive
    \item Plug the USB drive into your Raspberry Pi
    \item On the Pi, open the file manager and find your USB drive
    \item Copy the project folder to: \texttt{/home/pi/gesture\_recognition/}
    \item Or use terminal:
    \begin{lstlisting}
    cp -r /media/pi/USB_NAME/YourProjectFolder ~/gesture_recognition/
    \end{lstlisting}
    (Replace \texttt{USB\_NAME} with your actual USB drive name)
\end{enumerate}

\subsection{Step 6.3: Method 2 - Using SCP (From Another Computer)}

If you have another computer on the same network:

\begin{enumerate}
    \item Find your Raspberry Pi's IP address:
    \begin{lstlisting}
    hostname -I
    \end{lstlisting}
    You'll see something like: \texttt{192.168.1.100}
    
    \item On your other computer (Windows/Mac/Linux), open terminal/command prompt and type:
    \begin{lstlisting}
    scp -r YourProjectFolder pi@192.168.1.100:/home/pi/gesture_recognition/
    \end{lstlisting}
    (Replace \texttt{192.168.1.100} with your Pi's actual IP)
    
    \item Enter password when asked (default is \texttt{raspberry})
\end{enumerate}

\subsection{Step 6.4: Method 3 - Using Git}

If your project is on GitHub:

\begin{lstlisting}
cd ~/gesture_recognition
git clone https://github.com/yourusername/yourproject.git
\end{lstlisting}

\subsection{Step 6.5: Verify Files Are There}

Check that all files are in place:

\begin{lstlisting}
cd ~/gesture_recognition
ls -la
\end{lstlisting}

You should see \texttt{app.py} and the \texttt{model/} and \texttt{utils/} folders.

\section{Step 7: Set Up the Camera}

\subsection{What is This Step?}

The camera is the "eyes" of our system. We need to make sure the computer can see through it!

\subsection{Step 7.1: For USB Webcam}

\begin{enumerate}
    \item Plug your USB webcam into the Raspberry Pi
    \item Add yourself to the video group (gives permission to use camera):
    \begin{lstlisting}
    sudo usermod -a -G video $USER
    \end{lstlisting}
    
    \item Log out and log back in (or restart) for changes to take effect
    
    \item Test the camera:
    \begin{lstlisting}
    python3 -c "import cv2; cap = cv2.VideoCapture(0); print('Camera OK!' if cap.isOpened() else 'Camera FAILED'); cap.release()"
    \end{lstlisting}
    
    If you see "Camera OK!", you're good to go!
\end{enumerate}

\subsection{Step 7.2: For Raspberry Pi Camera Module}

\begin{enumerate}
    \item Enable the camera:
    \begin{lstlisting}
    sudo raspi-config
    \end{lstlisting}
    
    \item Navigate with arrow keys:
    \begin{itemize}
        \item Go to: \texttt{Interface Options}
        \item Select: \texttt{Camera}
        \item Choose: \texttt{Enable}
        \item Press \texttt{Finish}
    \end{itemize}
    
    \item Reboot:
    \begin{lstlisting}
    sudo reboot
    \end{lstlisting}
    
    \item Test the camera:
    \begin{lstlisting}
    raspistill -o test.jpg
    \end{lstlisting}
    
    If a file \texttt{test.jpg} appears, your camera works!
\end{enumerate}

\begin{tipbox}
If you have multiple cameras, you might need to use \texttt{--device 1} instead of \texttt{--device 0} in the app.py command.
\end{tipbox}

\section{Step 8: Optimize app.py for Raspberry Pi}

\subsection{What is This Step?}

Raspberry Pi is smaller and less powerful than a regular computer, so we need to make the program lighter (like using a smaller image size).

\subsection{Step 8.1: Check Current Settings}

Let's look at what resolution app.py uses by default. The default is 960x540, which might be too high for Pi.

\subsection{Step 8.2: Run with Lower Resolution}

Instead of changing the code, we can tell the program to use a smaller resolution when we run it:

\begin{lstlisting}
python3 app.py --width 640 --height 480
\end{lstlisting}

Or even smaller for better performance:

\begin{lstlisting}
python3 app.py --width 320 --height 240
\end{lstlisting}

\begin{infobox}
\textbf{Resolution explained:}
\begin{itemize}
    \item 960x540 = High quality, slower
    \item 640x480 = Medium quality, medium speed (recommended)
    \item 320x240 = Lower quality, faster
\end{itemize}
\end{infobox}

\section{Step 9: Run the Program!}

\subsection{What is This Step?}

Finally! Time to see your hand gesture recognition in action!

\subsection{Step 9.1: Make Sure You're in the Right Place}

\begin{lstlisting}
cd ~/gesture_recognition
source venv_mediapipe_pi/bin/activate
\end{lstlisting}

You should see \texttt{(venv\_mediapipe\_pi)} in your terminal.

\subsection{Step 9.2: Run app.py}

\begin{lstlisting}
python3 app.py --width 640 --height 480 --device 0
\end{lstlisting}

\begin{infobox}
\textbf{What the command means:}
\begin{itemize}
    \item \texttt{python3 app.py} = Run the app.py program
    \item \texttt{--width 640} = Use 640 pixels wide
    \item \texttt{--height 480} = Use 480 pixels tall
    \item \texttt{--device 0} = Use camera number 0 (first camera)
\end{itemize}
\end{infobox}

\subsection{Step 9.3: What You Should See}

A window should open showing:
\begin{itemize}
    \item Your camera feed
    \item Your hand with lines drawn on it (showing MediaPipe detecting your hand)
    \item Text showing what gesture it recognizes
    \item FPS (frames per second) counter
\end{itemize}

\subsection{Step 9.4: Controls}

\begin{itemize}
    \item \textbf{ESC} = Exit the program
    \item \textbf{0-9} = Select gesture number (for training mode)
    \item \textbf{N} = Normal mode
    \item \textbf{K} = Logging Key Point mode
    \item \textbf{H} = Logging Point History mode
\end{itemize}

\section{Step 10: Create a Startup Script}

\subsection{What is This Step?}

Instead of typing all those commands every time, we'll create a shortcut script!

\subsection{Step 10.1: Create the Script}

Create a new file:

\begin{lstlisting}
nano run_app.sh
\end{lstlisting}

This opens a text editor. Type these lines:

\begin{lstlisting}
#!/bin/bash
cd ~/gesture_recognition
source venv_mediapipe_pi/bin/activate
python3 app.py --width 640 --height 480 --device 0
\end{lstlisting}

Press \texttt{Ctrl+X}, then \texttt{Y}, then \texttt{Enter} to save.

\subsection{Step 10.2: Make It Executable}

\begin{lstlisting}
chmod +x run_app.sh
\end{lstlisting}

\subsection{Step 10.3: Run the Script}

Now you can simply type:

\begin{lstlisting}
./run_app.sh
\end{lstlisting}

Much easier!

\section{Step 11: Daily Operations Guide}

\subsection{What is This Step?}

Now that everything is set up, here's how to use your Raspberry Pi every day to run the gesture recognition project.

\subsection{Step 11.1: Starting the Project (Every Time)}

Every time you want to run the project, follow these steps:

\textbf{Step 1: Open Terminal}
\begin{itemize}
    \item Desktop Mode: Click Terminal icon or press \texttt{Ctrl + Alt + T}
    \item SSH: Connect via SSH from your computer
\end{itemize}

\textbf{Step 2: Navigate to Project Folder}
\begin{lstlisting}
cd ~/gesture_recognition
\end{lstlisting}

\textbf{Step 3: Activate Virtual Environment}
\begin{lstlisting}
source venv_mediapipe_pi/bin/activate
\end{lstlisting}

You should see \texttt{(venv\_mediapipe\_pi)} appear in your prompt.

\textbf{Step 4: Run the App}
\begin{lstlisting}
python3 app.py --width 640 --height 480 --device 0
\end{lstlisting}

\begin{warningbox}
\textbf{Important:} You must activate the virtual environment every time you open a new terminal! If you don't see \texttt{(venv\_mediapipe\_pi)} in your prompt, the program won't work.
\end{warningbox}

\subsection{Step 11.2: Using the Quick Start Script}

Instead of typing all those commands, you can use the script we created:

\begin{lstlisting}
~/run_gesture.sh
\end{lstlisting}

Or if you saved it in the project folder:
\begin{lstlisting}
cd ~/gesture_recognition
./run_app.sh
\end{lstlisting}

\subsection{Step 11.3: Checking if Everything is Working}

Before running the app, you can check if everything is set up correctly:

\textbf{Check Python Version:}
\begin{lstlisting}
python3 --version
\end{lstlisting}
Should show Python 3.9, 3.10, or 3.11

\textbf{Check Installed Packages:}
\begin{lstlisting}
source venv_mediapipe_pi/bin/activate
pip list
\end{lstlisting}
You should see: numpy, opencv-python-headless, tensorflow-lite-runtime, mediapipe

\textbf{Check Camera:}
\begin{lstlisting}
python3 -c "import cv2; cap = cv2.VideoCapture(0); print('Camera OK!' if cap.isOpened() else 'Camera FAILED'); cap.release()"
\end{lstlisting}

\textbf{Check Project Files:}
\begin{lstlisting}
cd ~/gesture_recognition
ls -la                    # Check main files
ls -la model/             # Check model folder
ls -la utils/             # Check utils folder
\end{lstlisting}

\subsection{Step 11.4: Stopping the Program}

To stop the running program:
\begin{itemize}
    \item Press \texttt{ESC} in the camera window, OR
    \item Press \texttt{Ctrl+C} in the terminal
\end{itemize}

\subsection{Step 11.5: Restarting Your Raspberry Pi}

Sometimes you need to restart your Pi:

\textbf{From Terminal:}
\begin{lstlisting}
sudo reboot
\end{lstlisting}

\textbf{From Desktop:}
Menu → Shutdown → Reboot

\subsection{Step 11.6: Shutting Down Your Raspberry Pi}

\textbf{From Terminal:}
\begin{lstlisting}
sudo shutdown now
\end{lstlisting}

\textbf{From Desktop:}
Menu → Shutdown → Shutdown

\begin{warningbox}
\textbf{Important:} Always shut down properly! Don't just unplug the power. This can damage your SD card.
\end{warningbox}

\subsection{Step 11.7: Common Daily Operations}

\textbf{Viewing Files:}
\begin{lstlisting}
cat app.py              # View file contents
less app.py             # View with scroll (press 'q' to quit)
nano app.py             # Edit file in text editor
\end{lstlisting}

\textbf{Finding Your IP Address:}
\begin{lstlisting}
hostname -I
\end{lstlisting}

\textbf{Checking Disk Space:}
\begin{lstlisting}
df -h
\end{lstlisting}
Shows how much space is used/available

\textbf{Checking Memory Usage:}
\begin{lstlisting}
free -h
\end{lstlisting}
Shows RAM usage

\textbf{Checking Running Processes:}
\begin{lstlisting}
ps aux
\end{lstlisting}
Shows all running programs

\textbf{Killing a Stuck Program:}
\begin{lstlisting}
# First, find the process
ps aux | grep python

# Then kill it (replace PID with the number you see)
kill PID
\end{lstlisting}

\textbf{Checking Pi Temperature:}
\begin{lstlisting}
vcgencmd measure_temp
\end{lstlisting}
Good temperature is below 70°C

\section{Step 12: Auto-Start on Boot (Optional)}

\subsection{What is This Step?}

If you want the gesture recognition app to start automatically when your Raspberry Pi boots up, you can set up a systemd service. This is like setting an app to start when your computer turns on.

\subsection{Step 12.1: Create Systemd Service File}

Create a service file:
\begin{lstlisting}
sudo nano /etc/systemd/system/gesture-app.service
\end{lstlisting}

Type this content (copy exactly):
\begin{lstlisting}
[Unit]
Description=Hand Gesture Recognition App
After=network.target

[Service]
Type=simple
User=pi
WorkingDirectory=/home/pi/gesture_recognition
Environment="PATH=/home/pi/gesture_recognition/venv_mediapipe_pi/bin"
ExecStart=/home/pi/gesture_recognition/venv_mediapipe_pi/bin/python3 app.py --width 640 --height 480 --device 0
Restart=always
RestartSec=10

[Install]
WantedBy=multi-user.target
\end{lstlisting}

Save: Press \texttt{Ctrl+X}, then \texttt{Y}, then \texttt{Enter}

\subsection{Step 12.2: Enable and Start the Service}

\begin{lstlisting}
# Reload systemd to recognize the new service
sudo systemctl daemon-reload

# Enable the service (starts on boot)
sudo systemctl enable gesture-app

# Start the service now
sudo systemctl start gesture-app
\end{lstlisting}

\subsection{Step 12.3: Check Service Status}

\begin{lstlisting}
sudo systemctl status gesture-app
\end{lstlisting}

You should see "active (running)" in green.

\subsection{Step 12.4: Managing the Auto-Start Service}

\textbf{Stop the service:}
\begin{lstlisting}
sudo systemctl stop gesture-app
\end{lstlisting}

\textbf{Start the service:}
\begin{lstlisting}
sudo systemctl start gesture-app
\end{lstlisting}

\textbf{Restart the service:}
\begin{lstlisting}
sudo systemctl restart gesture-app
\end{lstlisting}

\textbf{Disable auto-start (but keep service):}
\begin{lstlisting}
sudo systemctl disable gesture-app
\end{lstlisting}

\textbf{View service logs:}
\begin{lstlisting}
sudo journalctl -u gesture-app -f
\end{lstlisting}
Press \texttt{Ctrl+C} to exit

\begin{tipbox}
If you enable auto-start, the app will run every time the Pi boots. Make sure your camera is connected before turning on the Pi!
\end{tipbox}

\section{Troubleshooting}

\subsection{Problem 1: MediaPipe Won't Install}

\textbf{Symptoms:} You see errors when trying to install MediaPipe.

\textbf{Solutions:}

\begin{enumerate}
    \item Try a specific version:
    \begin{lstlisting}
    pip install mediapipe==0.10.8
    \end{lstlisting}
    
    \item Make sure you have enough space:
    \begin{lstlisting}
    df -h
    \end{lstlisting}
    You need at least 2GB free.
    
    \item Increase swap space (temporary storage):
    \begin{lstlisting}
    sudo dphys-swapfile swapoff
    sudo nano /etc/dphys-swapfile
    \end{lstlisting}
    Change \texttt{CONF\_SWAPSIZE=100} to \texttt{CONF\_SWAPSIZE=2048}
    Then:
    \begin{lstlisting}
    sudo dphys-swapfile setup
    sudo dphys-swapfile swapon
    \end{lstlisting}
\end{enumerate}

\subsection{Problem 2: Camera Not Detected}

\textbf{Symptoms:} Program says it can't find the camera.

\textbf{Solutions:}

\begin{enumerate}
    \item Check if camera is detected:
    \begin{lstlisting}
    ls -l /dev/video*
    \end{lstlisting}
    
    \item Try a different camera number:
    \begin{lstlisting}
    python3 app.py --device 1
    \end{lstlisting}
    
    \item For USB camera, make sure you're in the video group:
    \begin{lstlisting}
    groups
    \end{lstlisting}
    You should see "video" in the list.
\end{enumerate}

\subsection{Problem 3: Program Runs Too Slowly}

\textbf{Symptoms:} Low FPS, laggy video.

\textbf{Solutions:}

\begin{enumerate}
    \item Use lower resolution:
    \begin{lstlisting}
    python3 app.py --width 320 --height 240
    \end{lstlisting}
    
    \item Lower detection confidence (faster but less accurate):
    \begin{lstlisting}
    python3 app.py --width 640 --height 480 --min_detection_confidence 0.5
    \end{lstlisting}
    
    \item Close other programs to free up memory
    
    \item Make sure nothing is blocking the camera's view
\end{enumerate}

\subsection{Problem 4: "Module Not Found" Error}

\textbf{Symptoms:} Python says it can't find a module (like cv2, mediapipe, etc.).

\textbf{Solutions:}

\begin{enumerate}
    \item Make sure virtual environment is activated:
    \begin{lstlisting}
    source venv_mediapipe_pi/bin/activate
    \end{lstlisting}
    You should see \texttt{(venv\_mediapipe\_pi)}.
    
    \item Reinstall the missing package:
    \begin{lstlisting}
    pip install package_name
    \end{lstlisting}
    (Replace \texttt{package\_name} with the missing module)
\end{enumerate}

\subsection{Problem 5: "Model File Not Found" Error}

\textbf{Symptoms:} Program says it can't find model files.

\textbf{Solutions:}

\begin{enumerate}
    \item Check you're in the right directory:
    \begin{lstlisting}
    pwd
    \end{lstlisting}
    Should show: \texttt{/home/pi/gesture\_recognition}
    
    \item Check model files exist:
    \begin{lstlisting}
    ls -la model/keypoint_classifier/
    ls -la model/point_history_classifier/
    \end{lstlisting}
    
    \item Make sure all files were transferred correctly
\end{enumerate}

\subsection{Problem 6: "Command Not Found" Error}

\textbf{Symptoms:} Terminal says a command doesn't exist.

\textbf{Solutions:}

\begin{enumerate}
    \item Check spelling - Linux is case-sensitive!
    \item Make sure you're in the right directory
    \item Check if the command is installed:
    \begin{lstlisting}
    which command_name
    \end{lstlisting}
    \item For Python commands, make sure virtual environment is activated
\end{enumerate}

\subsection{Problem 7: "Permission Denied" Error}

\textbf{Symptoms:} Terminal says you don't have permission to do something.

\textbf{Solutions:}

\begin{enumerate}
    \item Use \texttt{sudo} before the command:
    \begin{lstlisting}
    sudo command_name
    \end{lstlisting}
    \item Check file permissions:
    \begin{lstlisting}
    ls -la filename
    \end{lstlisting}
    \item Change file permissions if needed:
    \begin{lstlisting}
    chmod +x filename    # Make file executable
    \end{lstlisting}
\end{enumerate}

\subsection{Problem 8: Virtual Environment Not Activating}

\textbf{Symptoms:} \texttt{source venv\_mediapipe\_pi/bin/activate} doesn't work.

\textbf{Solutions:}

\begin{enumerate}
    \item Make sure you're in the right directory:
    \begin{lstlisting}
    cd ~/gesture_recognition
    \end{lstlisting}
    \item Check if virtual environment exists:
    \begin{lstlisting}
    ls -la venv_mediapipe_pi/
    \end{lstlisting}
    \item If it doesn't exist, recreate it:
    \begin{lstlisting}
    python3 -m venv venv_mediapipe_pi
    \end{lstlisting}
    \item Make sure you're using the full path:
    \begin{lstlisting}
    source ~/gesture_recognition/venv_mediapipe_pi/bin/activate
    \end{lstlisting}
\end{enumerate}

\subsection{Problem 9: Out of Memory}

\textbf{Symptoms:} Program crashes, system becomes slow, "out of memory" errors.

\textbf{Solutions:}

\begin{enumerate}
    \item Check memory usage:
    \begin{lstlisting}
    free -h
    \end{lstlisting}
    \item Close other programs
    \item Increase swap space:
    \begin{lstlisting}
    sudo dphys-swapfile swapoff
    sudo nano /etc/dphys-swapfile
    \end{lstlisting}
    Change \texttt{CONF\_SWAPSIZE=100} to \texttt{CONF\_SWAPSIZE=2048}
    \begin{lstlisting}
    sudo dphys-swapfile setup
    sudo dphys-swapfile swapon
    \end{lstlisting}
    \item Use lower resolution in app.py
\end{enumerate}

\subsection{Problem 10: Pi Overheating}

\textbf{Symptoms:} Pi becomes slow, temperature warning, system shuts down.

\textbf{Solutions:}

\begin{enumerate}
    \item Check temperature:
    \begin{lstlisting}
    vcgencmd measure_temp
    \end{lstlisting}
    \item Add a cooling fan or heatsink
    \item Reduce CPU usage (lower resolution, close other programs)
    \item Make sure Pi has good ventilation
    \item Check if Pi is overclocked (reduce if needed)
\end{enumerate}

\subsection{Problem 11: Can't See Camera Window (SSH)}

\textbf{Symptoms:} Running via SSH, program starts but no window appears.

\textbf{Solutions:}

\begin{enumerate}
    \item This is normal! SSH doesn't show GUI windows
    \item Use X11 forwarding (advanced):
    \begin{lstlisting}
    ssh -X pi@192.168.1.100
    \end{lstlisting}
    \item Or use VNC to see the desktop remotely
    \item Or connect a monitor directly to the Pi
\end{enumerate}

\subsection{Problem 12: Files Not Transferring Correctly}

\textbf{Symptoms:} Files missing, wrong file sizes, errors when running.

\textbf{Solutions:}

\begin{enumerate}
    \item Check file sizes match:
    \begin{lstlisting}
    ls -lh filename
    \end{lstlisting}
    \item Verify all required files are present:
    \begin{lstlisting}
    find ~/gesture_recognition -name "*.tflite"
    find ~/gesture_recognition -name "*.csv"
    \end{lstlisting}
    \item Re-transfer files if needed
    \item Check file permissions:
    \begin{lstlisting}
    chmod 644 filename    # Make readable
    \end{lstlisting}
\end{enumerate}

\section{Quick Reference Commands}

Here's a comprehensive cheat sheet of all the important commands you'll need:

\subsection{Project Navigation}

\begin{lstlisting}
# Go to project folder
cd ~/gesture_recognition

# Go to home folder
cd ~

# See where you are
pwd

# List files
ls
ls -la              # Show all files including hidden
\end{lstlisting}

\subsection{Virtual Environment}

\begin{lstlisting}
# Activate virtual environment
source venv_mediapipe_pi/bin/activate

# Deactivate virtual environment
deactivate

# Check if activated (should see (venv_mediapipe_pi) in prompt)
# If not visible, activate it!
\end{lstlisting}

\subsection{Running the App}

\begin{lstlisting}
# Standard run (medium quality, recommended)
python3 app.py --width 640 --height 480 --device 0

# Low quality, faster
python3 app.py --width 320 --height 240 --device 0

# High quality, slower
python3 app.py --width 960 --height 540 --device 0

# Using quick start script
~/run_gesture.sh
# or
./run_app.sh
\end{lstlisting}

\subsection{System Checks}

\begin{lstlisting}
# Check Python version
python3 --version

# Check camera
python3 -c "import cv2; cap = cv2.VideoCapture(0); print('OK' if cap.isOpened() else 'FAIL'); cap.release()"

# Check installed packages
pip list

# Check Pi temperature
vcgencmd measure_temp

# Check disk space
df -h

# Check memory
free -h

# Check IP address
hostname -I
\end{lstlisting}

\subsection{File Operations}

\begin{lstlisting}
# View file
cat filename
less filename          # Scrollable view (press 'q' to quit)

# Edit file
nano filename          # Text editor (Ctrl+X to save and exit)

# Copy file
cp source destination

# Move/rename file
mv oldname newname

# Delete file (be careful!)
rm filename
rm -r foldername       # Delete folder

# Make file executable
chmod +x filename
\end{lstlisting}

\subsection{System Management}

\begin{lstlisting}
# Update system
sudo apt update
sudo apt upgrade -y

# Restart Pi
sudo reboot

# Shut down Pi
sudo shutdown now

# Check running processes
ps aux

# Kill a process
kill PID               # Replace PID with process number
\end{lstlisting}

\subsection{Service Management (Auto-Start)}

\begin{lstlisting}
# Check service status
sudo systemctl status gesture-app

# Start service
sudo systemctl start gesture-app

# Stop service
sudo systemctl stop gesture-app

# Restart service
sudo systemctl restart gesture-app

# Enable auto-start
sudo systemctl enable gesture-app

# Disable auto-start
sudo systemctl disable gesture-app

# View service logs
sudo journalctl -u gesture-app -f
\end{lstlisting}

\subsection{Emergency Commands}

\begin{lstlisting}
# Stop running program
Ctrl+C                 # In terminal

# Exit program window
ESC                    # In camera window

# Force quit terminal
Ctrl+D

# Clear terminal screen
clear
\end{lstlisting}

\section{Daily Workflow Summary}

\subsection{Your Daily Routine}

Here's the simple workflow you'll follow every time you want to use the gesture recognition:

\textbf{Starting Up:}
\begin{enumerate}
    \item Turn on Raspberry Pi (or connect via SSH)
    \item Open Terminal
    \item Run these commands:
    \begin{lstlisting}
    cd ~/gesture_recognition
    source venv_mediapipe_pi/bin/activate
    python3 app.py --width 640 --height 480 --device 0
    \end{lstlisting}
    \item Or use the quick script:
    \begin{lstlisting}
    ~/run_gesture.sh
    \end{lstlisting}
\end{enumerate}

\textbf{Using the App:}
\begin{itemize}
    \item Position your hand in front of the camera
    \item Make gestures (open hand, closed fist, pointing, etc.)
    \item Watch the recognition on screen
    \item Press \texttt{ESC} to exit when done
\end{itemize}

\textbf{Shutting Down:}
\begin{enumerate}
    \item Exit the app (press \texttt{ESC})
    \item Optionally shut down the Pi:
    \begin{lstlisting}
    sudo shutdown now
    \end{lstlisting}
\end{enumerate}

\begin{tipbox}
\textbf{Pro Tip:} If you set up auto-start (Step 12), the app will start automatically when the Pi boots. Just turn on the Pi and wait for the camera window to appear!
\end{tipbox}

\section{Performance Tips}

\subsection{For Better Speed}

\begin{itemize}
    \item Use 320x240 resolution instead of 640x480
    \item Close other programs
    \item Use a USB 3.0 webcam if available
    \item Make sure Raspberry Pi has good cooling (add a fan if needed)
\end{itemize}

\subsection{For Better Accuracy}

\begin{itemize}
    \item Use 640x480 or higher resolution
    \item Ensure good lighting
    \item Keep hand clearly visible
    \item Use \texttt{--min\_detection\_confidence 0.7} (default)
\end{itemize}

\section{What's Next?}

Now that you have it working, you can:

\begin{itemize}
    \item Experiment with different gestures
    \item Modify the code to add new features
    \item Connect it to control other devices (like in test.py)
    \item Train your own gesture models
\end{itemize}

\section{Summary}

Congratulations! You've successfully set up hand gesture recognition on Raspberry Pi! Here's everything we covered:

\subsection{Setup Steps Completed}

\begin{enumerate}
    \item Learned two ways to use Raspberry Pi (Desktop and SSH)
    \item Learned essential Linux commands
    \item Set up Raspberry Pi for first time
    \item Updated system and installed dependencies
    \item Created virtual environment
    \item Installed Python packages (NumPy, OpenCV, TensorFlow Lite, MediaPipe)
    \item Transferred project files
    \item Set up the camera (USB or Pi Camera Module)
    \item Optimized app.py for Raspberry Pi
    \item Ran the program successfully
    \item Created startup scripts for easy access
    \item Learned daily operations and maintenance
    \item Set up auto-start (optional)
    \item Learned troubleshooting techniques
\end{enumerate}

\subsection{Key Things to Remember}

\begin{tipbox}
\textbf{Every time you open a new terminal, you need to:}
\begin{enumerate}
    \item \texttt{cd ~/gesture\_recognition}
    \item \texttt{source venv\_mediapipe\_pi/bin/activate}
    \item Then run your program!
\end{enumerate}
\end{tipbox}

\begin{warningbox}
\textbf{Important Reminders:}
\begin{itemize}
    \item Always activate the virtual environment before running the app
    \item Check that camera is connected before starting
    \item Use proper shutdown (\texttt{sudo shutdown now}) - don't just unplug!
    \item Keep your Pi cool - add a fan if it gets too hot
    \item Check disk space regularly (\texttt{df -h})
\end{itemize}
\end{warningbox}

\subsection{Quick Start Checklist}

Before running the app, verify:
\begin{itemize}
    \item[$\square$] Virtual environment is activated (see \texttt{(venv\_mediapipe\_pi)} in prompt)
    \item[$\square$] You're in the project folder (\texttt{cd ~/gesture\_recognition})
    \item[$\square$] Camera is connected and working
    \item[$\square$] All project files are present (\texttt{app.py}, \texttt{model/}, \texttt{utils/})
    \item[$\square$] Python packages are installed (\texttt{pip list})
\end{itemize}

\subsection{Next Steps}

Now that everything is working, you can:
\begin{itemize}
    \item Experiment with different gestures
    \item Modify the code to add new features
    \item Connect hardware control (like in \texttt{test.py})
    \item Train your own gesture models
    \item Set up the app to run automatically on boot
    \item Share your project with others
\end{itemize}

\section{Appendix: Complete Setup Script}

If you want to automate everything, here's a complete setup script. Save it as \texttt{setup.sh}:

\begin{lstlisting}
#!/bin/bash
echo "Setting up MediaPipe on Raspberry Pi..."

# Update system
sudo apt update && sudo apt upgrade -y

# Install dependencies
sudo apt install -y build-essential cmake python3-pip python3-venv
sudo apt install -y libopencv-dev python3-opencv libatlas-base-dev
sudo apt install -y libjpeg-dev libtiff5-dev libjasper-dev libpng12-dev
sudo apt install -y libavcodec-dev libavformat-dev libswscale-dev libv4l-dev

# Create project folder
mkdir -p ~/gesture_recognition
cd ~/gesture_recognition

# Create virtual environment
python3 -m venv venv_mediapipe_pi
source venv_mediapipe_pi/bin/activate

# Install packages
pip install --upgrade pip setuptools wheel
pip install numpy
pip install opencv-python-headless
pip install tensorflow-lite-runtime
pip install mediapipe

echo "Setup complete!"
echo "Don't forget to transfer your project files!"
echo "Then activate with: source venv_mediapipe_pi/bin/activate"
\end{lstlisting}

Make it executable and run:
\begin{lstlisting}
chmod +x setup.sh
./setup.sh
\end{lstlisting}

\end{document}